\section{分时系统把中断、地址转换和保存现场连成进程}

早期批处理系统一次把一个作业交给机器；作业等待慢速外设时，昂贵的 CPU 也
跟着空闲。多道程序设计让多个作业同时驻留内存，一个作业等待时运行另一个。
交互式分时又加入周期时钟，使不主动让出 CPU 的计算也会被操作系统收回。
这里出现了进程的三个根本组成：

\begin{enumerate}[label=\arabic*.]
  \item 地址空间回答“这条地址属于哪些物理页、具有什么权限”；
  \item 处理器现场回答“下一条指令和所有寄存器从什么状态继续”；
  \item 调度状态回答“谁有资格在何时取得处理器”。
\end{enumerate}

仅有 PID 表不是进程，仅轮流调用函数也不是抢占。v0.9 的验收要求 QEMU 中的
x86-64 CPU 真实响应 8254 PIT IRQ0，保存 CPL3 现场，切换 CR3 与
TSS.RSP0，再从另一份现场执行 \code{iretq}。QEMU 只产生硬件事件；上述
所有状态和策略均由自研代码维护。

\begin{pdfdiagram}
  \node[os hardware] (pit) {8254 PIT\\IRQ0};
  \node[os block,right=of pit] (entry) {IDT/汇编入口\\176 B 帧};
  \node[os state,right=of entry] (scheduler) {round-robin\\Ready/Running};
  \node[os block,below=of scheduler] (switch) {CR3 + TSS.RSP0\\帧地址};
  \node[os state,left=of switch] (process) {另一进程\\同址用户 ELF};
  \draw[os arrow] (pit) -- (entry);
  \draw[os arrow] (entry) -- (scheduler);
  \draw[os arrow] (scheduler) -- (switch);
  \draw[os arrow] (switch) -- (process);
\end{pdfdiagram}
\diagramnote{时钟只提供机制；选择下一个 Ready 进程是策略；页表和内核栈使策略结果可由硬件安全执行。}

\section{v0.9 只实现拥有完整语义的四态模型}

本阶段固定四个 PCB，状态为
\code{Unused -> Ready -> Running -> Terminated}。创建失败可以把尚未运行
的 Ready 槽回滚为 Unused；启动选择第一个 Ready；时间片到期把当前 Running
放回 Ready 并选择后继；退出或用户异常进入 Terminated。

\begin{pdfdiagram}
  \node[os state] (unused) {Unused};
  \node[os block,right=of unused] (ready) {Ready};
  \node[os state,right=of ready] (running) {Running};
  \node[os hardware,right=of running] (terminated) {Terminated};
  \draw[os arrow] (unused) -- node[above]{创建} (ready);
  \draw[os arrow] (ready) -- node[above]{派发} (running);
  \draw[os arrow] (running) to[bend left] node[above]{抢占} (ready);
  \draw[os arrow] (running) -- node[above]{退出/异常} (terminated);
  \draw[os arrow] (ready) to[bend right] node[below]{创建回滚} (unused);
\end{pdfdiagram}
\diagramnote{本阶段没有等待条件，所以不伪造 Blocked；也没有父进程 wait，所以不伪造 Zombie。}

PCB 把地址空间、保存帧和可观测结果聚合：

\begin{longtable}{L{0.25\textwidth}L{0.27\textwidth}L{0.36\textwidth}}
\toprule
字段组 & 代表字段 & 所有权/不变量\\
\midrule
标识与状态 & 64 位 PID、状态、槽索引 & PID 单调；恰有零或一个 Running\\
地址空间 & PML4 根、入口、用户栈顶、页数 & 根不等于永久内核根；退出后释放\\
执行现场 & Ring 0 栈内的帧指针 & 完整 176 B 落在本 PCB 栈范围\\
终止结果 & exit code、异常向量/error/RIP/CR2 & 只在对应终止原因下有效\\
统计 & syscall、run tick、dispatch & 只在短中断临界路径更新\\
\bottomrule
\end{longtable}

PID 从 1 开始递增。创建 ELF 失败会释放 Ready 槽，但不会倒退 PID；已经写入
日志或被用户观察的标识不能悄悄指代另一个对象。固定容量不是通用限制，而是
当前教学增量：它让 PCB 本身不依赖尚无通用释放能力的早期堆。

\section{相同虚拟地址通过不同 CR3 到达不同物理页}

x86-64 四级页表把 48 位规范虚拟地址拆成四个 9 位索引和 12 位页内偏移：
\[
  v =
  \mathrm{PML4}[47{:}39]\,
  \mathrm{PDPT}[38{:}30]\,
  \mathrm{PD}[29{:}21]\,
  \mathrm{PT}[20{:}12]\,
  \mathrm{offset}[11{:}0].
\]
用户程序基址 \code{0x40000000} 的 PML4 索引为 0、PDPT 索引为 1；用户
栈接近低半规范地址顶端，PML4 索引为 255。这个布局允许进程只克隆一小段
页表，而不复制所有内核映射。

\begin{pdfdiagram}
  \node[os state] (template) {永久内核 PML4};
  \node[os block,below left=of template] (p1) {PID 1 PML4\\独立根};
  \node[os block,below right=of template] (p2) {PID 2 PML4\\独立根};
  \node[os state,below=of p1] (pdpt1) {克隆 PDPT[0]\\[0] 共享内核\\[1] 独占程序};
  \node[os state,below=of p2] (pdpt2) {克隆 PDPT[0]\\[0] 共享内核\\[1] 独占程序};
  \node[os hardware,below=of pdpt1] (leaf1) {VA 0x40000000\\物理帧 A};
  \node[os hardware,below=of pdpt2] (leaf2) {VA 0x40000000\\物理帧 B};
  \draw[os arrow] (template) -- (p1);
  \draw[os arrow] (template) -- (p2);
  \draw[os arrow] (p1) -- (pdpt1);
  \draw[os arrow] (p2) -- (pdpt2);
  \draw[os arrow] (pdpt1) -- (leaf1);
  \draw[os arrow] (pdpt2) -- (leaf2);
\end{pdfdiagram}
\diagramnote{虚拟地址只有和 CR3 选择的根放在一起才有意义；同址不是共享物理页的证据。}

创建地址空间时，页表管理器分配新 PML4 和新低端 PDPT，复制模板 PML4，
清空 PML4[255]；再复制模板低端 PDPT，清空 PDPT[1]。PDPT[0] 仍指向
共享 supervisor 内核子树。随后映射用户程序时，\code{MapPage} 为
PDPT[1] 分配独立 PD、PT 和叶帧；映射用户栈时从空 PML4[255] 分配另一
棵独占树。

PML4[0] 会因程序映射而带 U/S=1，但共享内核路径的 PDPT[0] 仍为
U/S=0。x86 权限取路径各级限制的交集，所以 CPL3 仍不能访问低端内核。
这说明“某一级 user”既不自动授权所有后代，也不能省略叶项和其他中间项检查。

写 CR3 会使当前处理器的大多数非 global TLB 翻译失效。PCID 可以减少这种
代价，但会新增标识分配、复用和定向失效协议。v0.9 先接受完整刷新，优先让
所有权和故障边界可解释。

\section{每进程 Ring 0 栈保存可跨时间片存活的现场}

从 CPL3 接受中断时，CPU 从 TSS.RSP0 取得新栈，依次保存旧 SS、RSP、
RFLAGS、CS 和 RIP。IRQ 桩再补 vector/error，并保存 15 个通用寄存器。
内存从低地址到高地址形成：

\begin{longtable}{L{0.22\textwidth}L{0.20\textwidth}L{0.46\textwidth}}
\toprule
偏移范围 & 长度 & 内容\\
\midrule
\code{0..119} & 120 B & R15..RAX 共 15 个 64 位通用寄存器\\
\code{120..135} & 16 B & vector 与规范化 error code\\
\code{136..159} & 24 B & RIP、CS、RFLAGS\\
\code{160..175} & 16 B & 旧用户 RSP、SS\\
\bottomrule
\end{longtable}

普通 Ring 0 异常没有特权切换，只有前 160 B；C++ 用保存 CS 的 RPL 判断
是否可以把公共帧扩展解释为 176 B 用户帧。首次运行没有硬件生成的旧现场，
创建代码在进程栈顶下方人工构造同一布局；因此首次进入和后续恢复共用一条
汇编弹栈路径。

每个槽位拥有 20 KiB 连续存储：底部 4 KiB 保持 not-present，其上 16 KiB
可用。guard 不是数组边界检查；它依赖当前页表真的省略叶映射。内存管理器在
建立低端身份映射时遍历四个 guard 地址并跳过它们，宿主布局测试则验证对齐、
相邻关系和帧范围。

若多个进程共用一块 RSP0 栈，第二个进程的 IRQ 或系统调用会覆盖第一进程挂起
的现场。切换 CR3 后必须同步把 TSS.RSP0 改为新栈顶；TSS 是每 CPU 状态，
不是自动随 CR3 切换的进程字段。

\section{汇编只搬运帧，C++ 只作有界决策}

异常、硬件 IRQ 和系统调用公共入口都遵循同一边界：

\begin{enumerate}[label=\arabic*.]
  \item 汇编清 DF、保存通用寄存器，把当前 RSP 作为首参数；
  \item 调整 RSP 到 System V AMD64 调用对齐并调用 C++；
  \item C++ 返回“下一份要恢复的帧地址”；
  \item 汇编用返回值替换 RSP，逆序恢复寄存器，丢弃 vector/error；
  \item \code{iretq} 恢复该帧中的处理器状态。
\end{enumerate}

不切换时 C++ 返回原帧；抢占、退出或用户异常时返回另一 PCB 的帧。汇编因此
不认识 Ready 队列、PID 或量子，C++ 也不复制十五个寄存器。边界的代价是
返回指针必须严格验证：当前帧应属于当前进程栈，目标帧应属于目标进程栈，
当前 CR3 应与当前 PCB 一致。

切换发生在 interrupt gate 清 IF 的窗口内。v0.9 单核环境不会在修改 PCB、
CR3 和 RSP0 的中间再次进入普通可屏蔽 IRQ；这不是多核锁。SMP 后另一 CPU
仍可并发访问共享队列，必须加入原子顺序和锁。

\section{四 tick 轮转用守恒性质证明公平}

PIT 目标约 1000 Hz，默认量子 \(q=4\) tick。每个用户 IRQ0 先增加当前
进程 run tick 和全局 timer tick，再增加已用预算。若预算小于 \(q\)，原帧
返回；预算达到 \(q\) 时，从当前槽位之后循环寻找 Ready：
\[
  \mathrm{candidate}(k) =
  (\mathrm{current}+1+k)\bmod 4,\qquad k=0,1,2,3.
\]
找不到其他 Ready 时只重置预算，不计抢占；找到时旧 Running 变 Ready，
候选变 Running，并同时增加 preemption 与 dispatch。

三个进程、总计 24 tick、量子 4 的纯模型集成测试中：
\[
  \mathrm{preemptions}=24/4=6,\qquad
  \mathrm{runTicks}_i=24/3=8.
\]
固定种子随机测试再生成 4096 组进程数、量子和 tick 序列。每一步检查：
\[
  \sum_i [\mathrm{state}_i=\mathrm{Running}] = 1,\qquad
  \sum_i \mathrm{runTicks}_i = \mathrm{timerTicks}.
\]
这些守恒比断言某一长序列更强：它们覆盖不同量子和队列长度，又不会把实现
内部循环写死成测试。

\section{退出路径先离开将被释放的地址空间}

当前进程调用 \code{ExitProcess} 或发生用户异常后，运行时记录终止原因并
让调度器选择后继。此时 C++ 仍在当前进程 Ring 0 栈上，当前 CR3 也仍指向
将被销毁的 PML4。若立即释放根，下一次取指或栈访问会使用已归还帧。

正确顺序为：

\begin{enumerate}[label=\arabic*.]
  \item 保存当前帧和终止现场，状态改为 Terminated；
  \item 得到后继索引或“全部完成”决定；
  \item 激活永久内核 CR3；
  \item 递归释放程序叶页、用户栈页、中间表、克隆 PDPT 和 PML4；
  \item 有后继则激活其 CR3/RSP0 并返回其帧；
  \item 无后继则恢复默认 RSP0 和调度启动前的内核调用链。
\end{enumerate}

销毁函数只遍历已知独占的 PDPT[1] 与 PML4[255]。复制来的其他 PML4 项和
PDPT[0] 是共享内核所有权，绝不能递归释放。接口还拒绝销毁当前活动 CR3 和
永久内核根，把正确顺序从注释提升成运行时前置条件。

\section{三个同址 worker 同时证明抢占与地址隔离}

正常镜像创建 PID1 smoke，以及 PID2、PID3、PID4 三个完全相同的 worker
ELF。每个 worker 先调用 \code{GetProcessId}，再执行三轮各 1,500,000 次
递增。全局 \code{workerCounter} 位于每个地址空间相同 BSS 虚拟地址；第
\(r\) 轮结束必须满足：
\[
  \mathrm{workerCounter}=r\times1\,500\,000.
\]
若三个 CR3 错误共享叶帧，后运行实例看到的初值不为零，至少一个检查失败。
三个实例各输出一次 \code{ADDRESS\_SPACE\_ISOLATED}，同时长计算区间让
真实 IRQ0 有机会产生多次抢占。

QEMU 一次代表执行观察到 107 个调度 tick、24 次抢占和 28 次派发。具体数字
受 TCG 速度和 IRQ 指令边界影响，不能硬编码；稳定验收是创建/终止恰为 4、
抢占至少 1、派发至少 4、九条 PID 进度日志各一次、四份退出码为零。

IRQ0 热路径完全不写串口。全部进程结束后，内核才批量输出调度总量和每进程
run tick、dispatch、syscall。创建前与全部销毁后的物理页统计必须一致，
才能输出 \code{PROCESS\_RESOURCES\_RECLAIMED}。这避免日志反过来延长量子，
也避免把 Terminated 状态误当作内存已经释放。

\begin{contractbox}
调度正确性由三层证据共同成立：宿主纯模型证明状态与守恒，ELF 审计证明同址
程序契约，QEMU 证明真实 IRQ、CR3、TSS、IRETQ 和页回收。任一层不能替代
另外两层。
\end{contractbox}

\section{完整进程生命周期是一组受约束状态转换}

新进程从 created 进入 runnable，调度后成为 running；等待 I/O 或同步条件时
进入 blocked，条件满足后回到 runnable；退出后进入 zombie，父进程取得退出
状态并回收后才变为 reaped。每次转换由持有相应队列锁的代码完成。

\begin{pdfdiagram}
  \node[os state] (created) {created};
  \node[os block,right=of created] (runnable) {runnable};
  \node[os state,right=of runnable] (running) {running};
  \node[os block,below=of running] (blocked) {blocked};
  \node[os state,left=of blocked] (zombie) {zombie};
  \node[os hardware,left=of zombie] (reaped) {reaped};
  \draw[os arrow] (created) -- (runnable);
  \draw[os arrow] (runnable) -- node[above]{调度} (running);
  \draw[os arrow] (running) to[bend left] node[above]{抢占} (runnable);
  \draw[os arrow] (running) -- node[right]{等待} (blocked);
  \draw[os arrow] (blocked) -- node[below]{唤醒} (runnable);
  \draw[os arrow] (running) -- node[below]{退出} (zombie);
  \draw[os arrow] (zombie) -- node[below]{父进程回收} (reaped);
\end{pdfdiagram}
\diagramnote{线程状态与所在队列必须一致；一个线程不能同时位于运行队列和等待队列。}

进程对象与线程对象生命周期不同。最后一个线程退出使进程进入退出流程，但打开
文件、子进程关系和退出码可能继续存在到父进程 wait。引用计数与父子关系必须
防止过早释放。

\section{ELF 用户程序装载建立新的地址空间}

内核为进程创建空页表，映射用户代码、只读数据、可写数据、BSS 与用户栈，
权限来自 PT\_LOAD 且执行 W\^X。装载过程使用内核临时映射写入目标物理页，
不能让用户在段尚未完成时运行。

用户栈顶按 ABI 对齐，初始内容包含参数计数、参数指针、环境和辅助向量。即使
第一版只传空参数，也要保留版本化栈布局，避免以后修改入口契约。

\section{上下文切换把栈本身作为保存容器}

线程被调度前已经在内核栈上。切换例程压入被调用者保存寄存器，把当前 RSP
写入线程对象，读取下一线程 RSP，再按相反顺序返回。函数的 \code{ret} 实际
返回到下一线程先前暂停的位置。

首次运行线程没有历史调用栈，创建代码人工构造一个与切换例程期望一致的栈帧，
使第一次恢复落到 trampoline。trampoline 调用线程入口，返回后统一执行线程
退出，避免控制流落入未定义地址。

\section{轮转调度的正确性先于复杂策略}

每个可运行线程位于一个运行队列，当前线程不重复入队。时钟 tick 减少预算，
耗尽时设置 need-reschedule；真正切换发生在中断退出或安全调度点，避免在任意
持锁位置切走。

公平性可以用每线程运行 tick、最长等待和上下文切换记录验证。优先级加入后要
处理饥饿与优先级反转；多核加入后要处理本地队列、迁移和缓存亲和性。先证明
单核轮转状态机，再逐层增加策略。

\section{等待队列把条件、锁与线程状态绑在一起}

线程在持有条件锁时检查资源；不满足就加入等待队列，并以不可分割顺序释放锁、
标记 blocked 和调用调度器。唤醒方持同一锁修改条件，把一个或多个等待者移到
运行队列。

条件必须循环检查，因为虚假唤醒、广播和竞争都可能使资源在获得 CPU 前再次
消失。超时和信号中断让等待返回多个结果，API 用枚举区分 satisfied、timed
out、interrupted 和 object closed。

\section{信号量、互斥锁和条件变量解决不同问题}

信号量计数可用资源，down 在计数为零时等待，up 增加并唤醒；互斥锁表达唯一
所有者并支持优先级继承；条件变量不保存事件，只允许线程等待某个受锁保护的
谓词变化。把三者都实现成一个整数会丢失所有权和错误检查。

自旋锁不能跨睡眠持有，互斥锁不能在 IRQ 使用，条件变量必须配合可睡眠锁。
接口类型应使无效组合难以表达。

\section{管道把有限缓冲、端点和唤醒组成 IPC}

管道拥有环形字节缓冲、读端计数、写端计数和两组等待队列。空管道且仍有写端
时读者阻塞；空且写端全关时读返回 EOF；满管道时写者阻塞；读端全关时写返回
broken pipe。

小于协议原子上限的写要么整体进入缓冲，要么等待，避免多个写者字节交错。
关闭端点必须在持锁状态更新计数并唤醒对侧，否则阻塞线程可能永久等待已不可能
发生的事件。

\section{文件描述符表提供进程局部命名}

fd 是进程表索引，槽位引用 open file description。后者保存 inode/对象引用、
当前偏移、状态标志和引用计数。\code{dup} 新 fd 共享同一 description，
所以文件偏移也共享；再次 \code{open} 创建独立 description。

查找 fd 时先在表锁下取得稳定引用，再释放表锁执行可能阻塞的 I/O。若直接保存
裸指针，另一线程 close 可能在操作中释放对象。

\section{VFS 用对象操作统一接口但保留语义差异}

vnode 或 inode 接口包含 lookup、read、write、truncate 和 stat 等操作。
普通文件按偏移访问，目录只允许目录操作，设备节点把请求转给驱动，管道没有
持久 inode。统一分发不表示所有操作都必须成功。

路径解析从根或 cwd 开始，逐段处理分隔符、点和父目录，检查执行/搜索权限，
限制符号链接深度。每次跨目录都取得稳定引用，避免并发删除后使用悬空对象。

\section{磁盘格式采用显式编码而非 C++ 内存布局}

超级块记录 magic、版本、块大小、总块数和各区域位置；位图记录 inode 与数据
块所有权；inode 记录类型、长度和块指针；目录项记录 inode 号、类型、名称
长度和名称字节。

编码器用小端读写函数逐字段访问，并检查字段范围与块边界。结构体 padding、
编译器 ABI 和宿主端序都不能进入磁盘格式。格式版本变化通过新解码路径处理，
不能让旧内核误读新字段。

\section{块映射从小文件逐步扩展}

inode 可以先保存若干直接块指针，小文件无需额外元数据；超过后使用一级、
二级间接块。文件偏移先除以块大小得到逻辑块号，再沿映射找到物理块，页内余数
确定字节位置。

扩展文件要先分配块、初始化并更新映射，任一步失败都回滚尚未提交的所有权。
稀疏文件未分配的逻辑块读取为零，不能访问物理块零。

\section{缓冲缓存是带状态的共享对象表}

缓存键通常为设备与块号。首次请求创建 loading 条目，只有一个线程发起 I/O，
其他线程等待；完成后状态变 clean，修改后变 dirty，回写中变 writeback，
错误状态保存诊断。引用计数防止正在使用的块被淘汰。

淘汰策略只从无引用且非 writeback 的对象选择。回写失败不能把 dirty 静默
清除，调用者与同步操作需要观察错误。锁顺序避免在持全局哈希锁时等待磁盘。

\section{崩溃一致性关注每个持久化边界}

内存中的多步更新在正常执行时可以受锁保护，掉电却会在任意磁盘写之间截断。
仅调整写入顺序可以保证某些简单操作，但复杂元数据需要日志或写时复制。

预写日志事务先写描述与新块内容，刷新后写提交记录，再把更新复制到正式位置。
恢复扫描完整提交事务并重放，忽略没有提交记录的尾部。日志头、校验、环绕和
序列号防止把旧残留误认为新事务。

\section{fsync 的承诺必须跨越设备缓存}

写系统调用返回可能只表示数据进入页缓存。块请求完成也可能只表示进入磁盘易失
缓存。要求持久化时，文件系统先写数据与元数据，再发送 flush/barrier 并等待
设备确认。QEMU 缓存模式会影响测试语义，运行参数必须记录。

第一版文件系统可以采用同步写简化，但文档必须明确返回时的保证。以后增加缓存
不能悄悄削弱已公开语义。

\section{Shell 的一条管道命令穿过全部内核层}

执行 \code{producer | consumer} 时，Shell 创建管道和两个进程，把生产者
stdout 与消费者 stdin 复制到管道端点，关闭多余 fd，分别装载 ELF，再等待
子进程。数据流经过用户缓冲、系统调用、管道等待队列和调度器。

\begin{enumerate}[label=\arabic*.]
  \item 终端驱动向 Shell 提交一行输入。
  \item 解析器形成命令与重定向结构，不在内核中解析文本语法。
  \item Shell 通过系统调用创建管道和进程。
  \item 子进程地址空间装载用户 ELF 与初始栈。
  \item 调度器交替运行生产者与消费者。
  \item 管道满/空触发睡眠和唤醒。
  \item 退出状态经 wait 返回 Shell。
\end{enumerate}

这条路径适合整机验收，因为它连接了用户边界、进程、同步、VFS 和驱动；每一
层仍要有独立测试，避免端到端失败无法定位。

\section{资源回收按所有权图逆序展开}

进程退出先阻止新线程和系统调用，关闭 fd 并唤醒等待者，释放用户映射与物理页，
保留退出码和最小进程对象供父进程回收。持锁、活动 I/O 和共享对象引用使回收
不能简单递归 delete。

测试循环创建、执行和回收进程，比较前后物理页、内核堆对象与打开文件计数。
单次成功无法发现缓慢泄漏，压力回归才验证生命周期闭环。
